\documentclass{article}
\usepackage[utf8]{inputenc}
\usepackage{amsmath,amssymb}
\usepackage[most]{tcolorbox}
\usepackage{geometry}
\geometry{margin=2.5cm}

\newcommand{\step}[1]{\par\noindent\hangindent=3.5em\hangafter=1%
  \makebox[3.5em][l]{\textbf{#1.}}}
\newcommand{\steptag}[1]{\hfill{\small\textsf{[#1]}}}

\newtcolorbox{proofbox}[1]{
  colback=gray!5, colframe=gray!60,
  fonttitle=\bfseries, title={#1},
  breakable, enhanced,
  left=4pt, right=4pt, top=4pt, bottom=4pt,
  before skip=10pt, after skip=10pt
}

\begin{document}

\begin{proofbox}{After first fixing a particular universal e\_sim>0 witness...}
\step{1} After first fixing a particular universal e\_sim>0 witness furnished by lem-maincb-corner-equivalence and a particular universal e\_full>0 witness furnished by lem-maincb-full-corner-identification, fix one def-maincb-witness-ledger datum W whose existence is furnished by lem-maincb-reset-constant-ledger instantiated with those same e\_sim,e\_full witnesses; then 0 <= epsilon <= W.epsilon\_MAIN implies epsilon <= W.e\_env, epsilon <= W.e1/W.K1, epsilon <= W.e\_s2/W.K2, and epsilon <= W.e\_cross/W.K3, while the global scalar scale epsilon, atomic scalar scale W.K\_call*epsilon, and Stage-1, Stage-2, and Stage-3 scales W.K1*epsilon,W.K2*epsilon,W.K3*epsilon are all at most W.K\_call*epsilon <= W.r\_reset,e\_sim,e\_full; moreover W.L*epsilon <= W.K\_call*epsilon and W.c0\_cb*W.K\_call*epsilon <= 1/2. \steptag{validated} \\[6pt]
\step{1.1} With the particular e\_sim and e\_full witnesses fixed exactly as in the statement, lem-maincb-reset-constant-ledger, instantiated with those same witnesses and its compatible fixed provider witnesses, supplies a def-maincb-witness-ledger datum W all of whose fields are positive and for which W.K\_call=max\{1,W.L+1,W.c0\_cb,W.K1,W.K2,W.K3\} and W.epsilon\_MAIN=min\{W.e\_env,W.e1/W.K1,W.e\_s2/W.K2,W.e\_cross/W.K3,W.r\_reset/W.K\_call,e\_sim/W.K\_call,e\_full/W.K\_call,[2*max\{1,W.c0\_cb*W.K\_call\}]\textasciicircum{}(-1)\}. \steptag{validated} \\[6pt]
\step{1.1.1} By lem-maincb-corner-equivalence choose the particular universal e\_sim>0 witness, and by lem-maincb-full-corner-identification choose the particular universal e\_full>0 witness. Fix the compatible witnesses furnished by lem-maincb-error-improvement, lem-maincb-direct-corner-envelope, lem-maincb-stage1-call-envelope, lem-maincb-stage2-call-envelope, and lem-maincb-stage3-call-envelope in the order required by lem-maincb-reset-constant-ledger. Universal instantiation of lem-maincb-reset-constant-ledger with these exact e\_sim,e\_full witnesses then supplies the asserted datum W, explicitly described there as supplied by lem-maincb-witness-arithmetic. \steptag{validated} \\[6pt]
\step{1.1.1.1} \textsc{assume}: Scope correction: node 1.1.1 is not an unconditional derivation of the provider witnesses or of W from the five abbreviated dependencies named in its second sentence. The root statement begins only after (i) the particular e\_sim and e\_full witnesses have been fixed and (ii) one W has been fixed whose provenance is a complete instantiation of lem-maincb-reset-constant-ledger with those same witnesses. Consequently the standing context includes, inside that provenance condition, every antecedent witness explicitly enumerated in the imported statement of lem-maincb-reset-constant-ledger: the improvement-iteration witnesses, reset-invariant-preservation witnesses, initial-raw-inclusion witnesses, Stage-1 corner-unit and raw-refinement witnesses, Stage-2 extcb-datum and raw-extension witnesses, Stage-3 merging-datum and raw-merge witnesses, and the remaining listed compatible choices. No existence claim for any omitted provider is made in this workspace, and the five-lemma list in node 1.1.1 must not be read as exhaustive or as furnishing the omitted providers. \steptag{validated} \\[4pt]
\step{1.1.1.2} \textsc{q.e.d.}: Under that exact standing fixed-provenance context, apply conclusion-elimination to the already-complete instantiation recorded by the root phrase that W is furnished by lem-maincb-reset-constant-ledger instantiated with those same e\_sim,e\_full witnesses. The imported conclusion then identifies this already-fixed W as a def-maincb-witness-ledger datum supplied by lem-maincb-witness-arithmetic. This is the asserted datum used by node 1.1; it is a conditional use of the root fixed-W binder, not an attempted discharge of the provider antecedents from the allowed externals. \steptag{validated} \\[4pt]
\step{1.1.2} Unpacking the phrase supplied by lem-maincb-witness-arithmetic in the established conclusion of lem-maincb-reset-constant-ledger gives positivity of every field and the exact defining formulas W.K\_call=max\{1,W.L+1,W.c0\_cb,W.K1,W.K2,W.K3\} and W.epsilon\_MAIN=min\{W.e\_env,W.e1/W.K1,W.e\_s2/W.K2,W.e\_cross/W.K3,W.r\_reset/W.K\_call,e\_sim/W.K\_call,e\_full/W.K\_call,[2*max\{1,W.c0\_cb*W.K\_call\}]\textasciicircum{}(-1)\} for this same W and these same e\_sim,e\_full witnesses. \steptag{validated} \\[4pt]
\step{1.2} Assume 0<=epsilon<=W.epsilon\_MAIN. Since a finite minimum is at most each of its entries, the displayed formula for W.epsilon\_MAIN gives epsilon<=W.e\_env, epsilon<=W.e1/W.K1, epsilon<=W.e\_s2/W.K2, and epsilon<=W.e\_cross/W.K3. \steptag{validated} \\[4pt]
\step{1.3} The displayed maximum formula gives W.K\_call>=1,W.K1,W.K2,W.K3. Multiplication by epsilon>=0 therefore gives epsilon<=W.K\_call*epsilon and W.Ki*epsilon<=W.K\_call*epsilon for i=1,2,3; the atomic scalar scale equals W.K\_call*epsilon. Thus the global, atomic, and three stage scales are all at most W.K\_call*epsilon. \steptag{validated} \\[4pt]
\step{1.4} From the displayed minimum formula, epsilon<=W.epsilon\_MAIN is at most W.r\_reset/W.K\_call, e\_sim/W.K\_call, and e\_full/W.K\_call. Since W.K\_call>0, multiplication by W.K\_call yields W.K\_call*epsilon<=W.r\_reset, W.K\_call*epsilon<=e\_sim, and W.K\_call*epsilon<=e\_full. \steptag{validated} \\[6pt]
\step{1.4.1} Use the already validated node 1.1.2, rather than the still-pending aggregate node 1.1: for this same W and these same fixed e\_sim,e\_full it establishes W.K\_call>0 and the exact formula W.epsilon\_MAIN=min\{W.e\_env,W.e1/W.K1,W.e\_s2/W.K2,W.e\_cross/W.K3,W.r\_reset/W.K\_call,e\_sim/W.K\_call,e\_full/W.K\_call,[2*max\{1,W.c0\_cb*W.K\_call\}]\textasciicircum{}(-1)\}. Hence the elementary finite-minimum property and the standing hypothesis epsilon<=W.epsilon\_MAIN give epsilon<=W.r\_reset/W.K\_call, epsilon<=e\_sim/W.K\_call, and epsilon<=e\_full/W.K\_call. Multiplying each inequality by the positive real W.K\_call preserves its direction, and W.K\_call*(a/W.K\_call)=a for a=W.r\_reset,e\_sim,e\_full because W.K\_call is nonzero. Therefore W.K\_call*epsilon<=W.r\_reset, W.K\_call*epsilon<=e\_sim, and W.K\_call*epsilon<=e\_full. \steptag{validated} \\[4pt]
\step{1.5} The maximum formula gives W.K\_call>=W.L+1>=W.L, hence W.L*epsilon<=W.K\_call*epsilon for epsilon>=0. Also epsilon<=[2*max\{1,W.c0\_cb*W.K\_call\}]\textasciicircum{}(-1); multiplying by the positive W.c0\_cb*W.K\_call and using W.c0\_cb*W.K\_call<=max\{1,W.c0\_cb*W.K\_call\} gives W.c0\_cb*W.K\_call*epsilon<=1/2. \steptag{validated} \\[6pt]
\step{1.5.1} Use validated node 1.1.2 directly, not pending aggregate node 1.1. It gives positivity of W.c0\_cb and W.K\_call and the exact formulas W.K\_call=max\{1,W.L+1,W.c0\_cb,W.K1,W.K2,W.K3\} and W.epsilon\_MAIN=min\{W.e\_env,W.e1/W.K1,W.e\_s2/W.K2,W.e\_cross/W.K3,W.r\_reset/W.K\_call,e\_sim/W.K\_call,e\_full/W.K\_call,[2*max\{1,W.c0\_cb*W.K\_call\}]\textasciicircum{}(-1)\}. Thus W.K\_call>=W.L+1>=W.L, and multiplication by the standing epsilon>=0 gives W.L*epsilon<=W.K\_call*epsilon. Put x=W.c0\_cb*W.K\_call and M=max\{1,x\}. Positivity from node 1.1.2 gives x>0 and M>=1>0. The minimum formula together with epsilon<=W.epsilon\_MAIN gives epsilon<=1/(2M). Multiplication by x>0 yields x*epsilon<=x/(2M), while x<=M and 2M>0 give x/(2M)<=M/(2M)=1/2. Hence W.c0\_cb*W.K\_call*epsilon<=1/2. \steptag{validated} \\[4pt]
\end{proofbox}

\end{document}
